\documentclass[11pt]{article}

\usepackage[letterpaper,margin=0.9in]{geometry}
\usepackage[T1]{fontenc}
\usepackage{lmodern}
\usepackage{microtype}
\usepackage{amsmath,amssymb,mathtools,bm}
\usepackage{booktabs,array,tabularx,longtable}
\usepackage{xcolor}
\usepackage{enumitem}
\usepackage{graphicx}
\usepackage{float}
\usepackage{tikz}
\usetikzlibrary{arrows.meta,calc,positioning,fit,backgrounds}
\usepackage{url}
\usepackage[hidelinks]{hyperref}
\usepackage[nameinlink,noabbrev]{cleveref}

\definecolor{navy}{HTML}{15324A}
\definecolor{blue}{HTML}{2E6F9E}
\definecolor{green}{HTML}{348A59}
\definecolor{orange}{HTML}{C77700}
\definecolor{red}{HTML}{B64949}
\definecolor{purple}{HTML}{76538A}
\definecolor{graytext}{HTML}{5E6872}
\definecolor{lightblue}{HTML}{EAF2F8}
\definecolor{lightgray}{HTML}{F2F4F5}
\definecolor{lightred}{HTML}{FBEEEE}

\hypersetup{
  pdftitle={A Clairvoyant Benchmark for Fleet Dispatch},
  pdfauthor={Fleet Policy Experimentation Framework},
  colorlinks=true, linkcolor=navy, citecolor=blue, urlcolor=blue
}

\setlist[itemize]{leftmargin=1.35em,itemsep=0.25em,topsep=0.4em}
\setlist[enumerate]{leftmargin=1.55em,itemsep=0.25em,topsep=0.4em}
\renewcommand{\arraystretch}{1.15}
\newcommand{\R}{\mathbb{R}}
\newcommand{\Z}{\mathbb{Z}}
\newcommand{\tail}{\operatorname{tail}}
\newcommand{\head}{\operatorname{head}}
\newcommand{\src}{\mathrm{src}}
\newcommand{\snk}{\mathrm{snk}}
\newcommand{\svc}{\mathrm{svc}}
\newcommand{\emp}{\mathrm{emp}}
\newcommand{\chg}{\mathrm{chg}}
\newcommand{\que}{\mathrm{que}}
\newcommand{\rep}{\mathrm{rep}}
\newcommand{\idl}{\mathrm{idle}}
\newcommand{\dep}{\mathrm{dep}}
\newcommand{\rel}{\mathrm{rel}}


\title{\vspace{-1.25em}\textbf{A Clairvoyant Benchmark for Fleet Dispatch}}
\author{}
\date{}

\begin{document}
\maketitle
\vspace{-1.4em}

\begin{abstract}
An electric mobility service must assign vehicles to arriving requests while managing their location,
availability and stored energy. A policy comparison estimates the effect of changing those operating
decisions, but it does not show how close either policy came to the best result available from the
same fleet state and realized demand. This note presents a framework for constructing a Clairvoyant
(perfect information) benchmark for that purpose. For each realized day, a mixed-integer program observes the
complete request sequence and maximizes the number of riders collected within the service threshold.
The perfect-information optimum, denoted by $N^*$, is not computed directly because the formulation does
not represent the departure time of every vehicle movement. The program is solved with conservative and
optimistic travel durations. Execution of the conservative plan establishes an attainable value
$L\leq N^*$. The optimistic value $U$ satisfies $N^*\leq U$ when the formulation covers every trajectory
permitted by the operating model; the restrictions listed in \cref{sec:conditions} make $U$ conditional.
An operating policy may exceed $L$, while $N^*$ remains the perfect-information optimum. Results are
not published here.
\end{abstract}

\section{Operational setting and benchmark purpose}

The operating model represents a battery-electric fleet operating within a closed service region. Each
request has an arrival time, an origin, a destination and a rider-specific tolerance for the quoted
pickup time. An assigned vehicle drives empty to the origin, collects the rider and completes the
passenger journey. The vehicle remains unavailable until that work is complete. A request may instead
be declined, rejected by the system or abandoned while waiting.

At regular invocation times, a policy can assign available vehicles, direct idle vehicles to another
zone and send vehicles to charge. These actions compete for the same fleet capacity. Repositioning may
place a vehicle closer to a later request, but the movement consumes time and energy and may leave the
origin zone uncovered. Charging restores energy while temporarily removing a vehicle from service.
Assignments made now therefore change the locations, energy levels and availability times presented to
later requests.

\paragraph{Timely service.}
The primary outcome is timely service. In the Brooklyn configuration, a request is timely when pickup
occurs no more than ten minutes after the request arrives. The corresponding rate divides timely pickups
by every request arriving in the reporting window, including declines, system rejections and
abandonments. The benchmark maximizes the count because the number of arrivals is fixed within a
realized instance.

\paragraph{Choice of ten-minute threshold.}
The ten-minute threshold is a scenario parameter rather than an estimate fitted from observed pickup
times. In the pooled Uber and Lyft records used for the Monday--Wednesday demand calculation, completed
Brooklyn trips have a mean request-to-pickup interval of 4.5 minutes. That sample excludes requests that
were never served and is therefore unsuitable for estimating rider patience or a service requirement.
The observed interval also reflects operations in physical space, whereas the zonal model represents
vehicle and rider locations through aggregate zones.

The Brooklyn inputs aggregate the borough into 20 decision zones. Under that representation, none of the
400 expected zone-pair travel times is five minutes or less; the minimum is 7.2 minutes. Even pickup
within the rider's zone is represented by approximately 9.5--10.7 minutes of travel. A five-minute
threshold would consequently measure the lower tail of the zonal travel approximation and would leave
little scope for a dispatch policy to affect the result. Ten minutes was selected as the shortest
threshold that clears the model's travel-time floor while remaining restrictive. This choice is a
modeling tradeoff: modeled timely service at ten minutes should not be compared directly with the
observed 4.5-minute mean as if the two quantities were measured under the same spatial representation
or sampling process.

A paired experiment determines whether one policy improves timely service relative to
another under matched random inputs. Similar policy results remain ambiguous: the policies may both
be close to the best attainable result, or both may leave substantial capacity unused. A reference
optimization is needed to distinguish those cases.

The Clairvoyant (perfect information) benchmark provides that reference. It is computed
retrospectively from the realized requests, the rider attributes and the fleet state at the start of
the reporting window. It can coordinate vehicle paths with knowledge of future requests, while an
operating policy uses only the information available when it acts. The benchmark measures operational
headroom. Its use of future requests precludes deployment as an online policy.

Each benchmark instance is conditional on one realized request sequence, one reporting window and the
fleet state at the start of that window. Its conclusions apply to the declared operating and input
models. Physical fleet performance, latent demand and demand outside the modeled service region are
outside the resulting claim.

\section{Inputs and interpretation}
\label{sec:study}

\subsection{Input data}

The Brooklyn inputs are fitted from New York City Taxi and Limousine Commission High Volume For Hire
Vehicle records for June and July 2025. The demand artifact contains 7,089,190 completed trips with both
endpoints in Brooklyn, represented by day type, origin, destination and request hour. The travel
artifact contains 7,857,311 completed passenger journeys, represented by origin, destination and hour.
Brooklyn's 61 taxi zones are aggregated into 20 contiguous decision zones.

The source records contain completed trips rather than all attempted requests. Demand levels are
therefore conditional on the fitted completed-trip distribution and a declared capture share. The
capture share and the fleet size are scenario parameters rather than estimates of market share or of
the fleet required to serve Brooklyn. The benchmark measures attainable performance within this
constructed operating environment.

Travel dispersion is disabled when a benchmark instance is built. A movement then takes the expected
duration in the matrix cell selected by its origin, destination and departure hour, and the benchmark
cannot improve its objective by selecting favorable travel-time draws.

\subsection{Outputs and interpretation}

For each policy and realized day, the inputs to the benchmark are the policy's timely pickup count
$N_p$ and the fleet state at the start of the reporting window. The benchmark combines that state with
the realized requests during the reporting window and reports two optimization values, $L$ and $U$.
The paired policy experiment estimates the effect of changing the online policy. The benchmark
comparison assesses how much timely service remained attainable once the starting state had been fixed.

For interpretation, $N_p$ records what the evaluated online policy achieved, and $L$ records what one
executable plan with future information achieved. $U$ is an optimistic comparison whose upper-bound
interpretation depends on the coverage conditions in \cref{sec:conditions}. The mathematical target
$N^*$ is the best timely pickup count over all trajectories permitted by the operating model from the
same initial state.

A policy can exceed $L$ because $L$ is the value of one feasible conservative plan. If $N_p>L$, the
policy itself supplies a stronger attainable value, so the lower benchmark becomes $\max\{L,N_p\}$.
When the upper coverage condition also holds, a policy's share of the optimum satisfies
\begin{equation}
  \frac{N_p}{U}
  \leq
  \frac{N_p}{N^*}
  \leq
  \frac{N_p}{\max\{L,N_p\}}.
  \label{eq:share}
\end{equation}

\section{Perfect-information objective}
\label{sec:objective}

Fix a realized day, a reporting window $[T_0,T_1)$ and the fleet state at $T_0$. Let $\Sigma$ denote
the trajectories that can be generated from that state by sequences of instructions permitted by the
dispatch rules. For trajectory $\sigma\in\Sigma$, let $N(\sigma)$ be the number of riders whose requests
arrive in the reporting window and whose pickups occur within the service threshold $\delta$. The
perfect-information optimum is
\begin{equation}
  N^* = \max_{\sigma\in\Sigma} N(\sigma).
  \label{eq:nstar}
\end{equation}

For an evaluated policy, let $N_p$ be its timely pickup count on the same realized day. The policy
trajectory belongs to $\Sigma$, so
\begin{equation}
  N_p \leq N^*.
  \label{eq:policy-optimum}
\end{equation}
This relation follows from the definition of $N^*$ and does not depend on the mixed-integer model.

The benchmark observes the request time, origin and destination of every request in the window; the
rider's assignment deadline and acceptance threshold; the expected travel duration and distance for
each movement; and every vehicle's location, next available time and stored energy at $T_0$. These
quantities make each benchmark instance deterministic.

\section{Why two optimization problems are solved}
\label{sec:bounds}

Travel duration depends on departure hour. Passenger journeys begin at modeled pickup times, so the
formulation selects their departure hour exactly. Other movements begin at times that the formulation
does not represent explicitly. These include the initial movement from a vehicle to its first rider,
the approach between consecutive riders, movements to and from the depot, and the within-zone movement
to a pickup. Their feasible departure intervals can cross an hour boundary inside the reporting window.

An exact time-dependent formulation would add a departure time and hour-selection variables to every
affected movement, together with constraints permitting a vehicle to wait before departure. The
formulation instead assigns one duration to each affected movement before solving. Let
$\pi^{\min}$ assign the shortest expected duration over the feasible departure hours and let
$\pi^{\max}$ assign the longest. Write $V(\pi)$ for the optimal objective under duration assignment
$\pi$.

\paragraph{Optimistic solve.}
\begin{equation}
  U = V(\pi^{\min}).
\end{equation}
The minimum durations relax the travel-time component of the problem. If every trajectory permitted by the operating model
is represented by the remaining formulation, each trajectory is feasible in the optimistic solve and
$N^*\leq U$. This coverage condition is material; the current exceptions are listed in
\cref{sec:conditions}.

\paragraph{Conservative solve.}
\begin{equation}
  L = V(\pi^{\max}).
\end{equation}
The maximum durations produce a conservative rider schedule. That schedule is executed under the same
operating rules. When execution reproduces its rider-level outcomes, it provides a realized trajectory
with objective $L$, establishing $L\leq N^*$.

When both conditions hold, the two values give
\begin{equation}
  L \leq N^* \leq U.
  \label{eq:bracket}
\end{equation}

The approach keeps the mixed-integer program small and quick to solve while measuring the
approximation introduced by time-dependent travel: the difference $U-L$ bounds the effect of fixing
each affected movement's duration before solving.

\section{Mixed-integer formulation}
\label{sec:formulation}

\subsection{Instance data}

Let $\mathcal R$ be the requests arriving in the reporting window and $\mathcal V$ the vehicles in the
fleet. For request $j$, the instance contains request time $a_j$, origin $o_j$, destination $d_j$,
assignment deadline $\bar D_j$, acceptance threshold $\theta_j$, and within-zone pickup duration
$\iota_j$. Its feasible pickup interval is
\begin{equation}
  \ell_j = a_j + \iota_j,
  \qquad
  u_j = \min\{\bar D_j + \theta_j,\ T_1-T_0\}.
  \label{eq:pickup-window}
\end{equation}
A request with $\ell_j>u_j$ cannot be collected within the modeled window and is removed.

For vehicle $v$, the instance contains its zone $z_v$, next available time $F_v$ and stored energy
$B_v$ at $T_0$. The energy required for a movement from $o$ to $d$ is $\epsilon(o,d)$. Battery
capacity is $\bar B$, and the reserve floor $\underline B$ is the largest energy required to reach the
depot from any zone in the instance.

\subsection{Time-dependent passenger journeys}

Pickup time $t_j$ determines the departure hour of the passenger journey. The pickup interval is split
at each hour boundary into segments $s=0,\ldots,m_j$. Segment $s$ has endpoints
$[\sigma_{js},\varepsilon_{js}]$ and passenger journey duration $\rho_{js}$. Binary $h_{js}$ selects
the applicable segment:
\begin{align}
  \textstyle\sum_{s\geq1} h_{js} &\leq 1, \\
  t_j &\geq \sigma_{j0} + \textstyle\sum_{s\geq1}(\sigma_{js}-\sigma_{j0})h_{js}, \\
  t_j &\leq \varepsilon_{j0} + \textstyle\sum_{s\geq1}(\varepsilon_{js}-\varepsilon_{j0})h_{js}, \\
  \rho_j(h) &= \rho_{j0} + \textstyle\sum_{s\geq1}(\rho_{js}-\rho_{j0})h_{js}.
  \label{eq:ride-duration}
\end{align}
Passenger journey duration is therefore determined by the hour containing the modeled pickup time.

\subsection{Network and decision variables}

A vehicle path contains three types of arc: an entry arc $(v,j)\in\mathcal A^E$ from a vehicle's
initial state to its first rider; a rider arc $(i,j)\in\mathcal A^R$ between consecutive riders; and a
depot arc $(i,j,k)\in\mathcal A^D$ that inserts a charging visit of $k$ fifteen-minute increments
between two riders. An arc is omitted when its estimated arrival time would cause the rider to reject
the quote.

\begin{center}
\begin{tabularx}{\textwidth}{@{}llX@{}}
\toprule
Variable & Domain & Interpretation \\
\midrule
$y_j$ & $\{0,1\}$ & request $j$ is collected \\
$w_j$ & $\{0,1\}$ & request $j$ is collected within $\delta$ \\
$t_j$ & $[\ell_j,u_j]$ & pickup time of request $j$ \\
$h_{js}$ & $\{0,1\}$ & pickup $j$ occurs in hour segment $s\geq1$ \\
$e_{vj}$, $z_{ij}$, $x_{ijk}$ & $\{0,1\}$ & entry, rider or depot arc is selected \\
$\mu_{ij}$, $\widetilde\mu_{ijk}$ & $\{0,1\}$ & vehicle waits in $o_j$ before request $j$ arrives \\
$b_j$ & $[\underline B,\bar B]$ & energy stored after delivering request $j$ \\
\bottomrule
\end{tabularx}
\end{center}

\subsection{Principal constraints}

Each vehicle enters at most one path. Each collected request has one predecessor and at most one
successor:
\begin{align}
  \textstyle\sum_j e_{vj} &\leq 1 && \forall v, \\
  \textstyle\sum_v e_{vj} + \sum_i z_{ij} + \sum_{i,k}x_{ijk} &= y_j && \forall j, \\
  \textstyle\sum_{j'} z_{jj'} + \sum_{j',k}x_{jj'k} &\leq y_j && \forall j, \\
  w_j &\leq y_j && \forall j.
\end{align}

For rider arcs, the pickup of $j$ follows the pickup and passenger journey of $i$ and the approach from
$d_i$ to $o_j$. Depot arcs add the two depot movements and the charging interval:
\begin{align}
  t_j &\geq t_i + \rho_i(h) + p_{ij} - M_{ij}(1-z_{ij}), \\
  t_j &\geq t_i + \rho_i(h) + q_i + k\kappa + r_j - M_{ijk}(1-x_{ijk}).
  \label{eq:sequence}
\end{align}
Here $p_{ij}$ is the direct approach duration, $q_i$ the duration to the depot, $r_j$ the duration from
the depot, and $\kappa=15$ minutes. Each big-$M$ constant is set to the smallest value that deactivates
its row when the corresponding arc is not selected.

A vehicle that reaches $o_j$ before the request may wait and then perform the within-zone movement.
Otherwise it begins the full approach after assignment. For $p_{ij}>\iota_j$:
\begin{align}
  t_j &\geq a_j+p_{ij}-(p_{ij}-\iota_j)\mu_{ij}-M(1-z_{ij}), \\
  \mu_{ij} &\leq z_{ij}, \\
  t_i+\rho_i(h)+p_{ij} &\leq a_j-\tau+M(1-\mu_{ij}),
  \label{eq:preposition}
\end{align}
where $\tau$ is the policy invocation interval. Corresponding constraints apply to entry and depot
arcs.

The vehicle must commit by the rider's assignment deadline, and a timely pickup occurs within
$\delta$ of the request:
\begin{align}
  t_j-p_{ij} &\leq \bar D_j+M(1-z_{ij}), \\
  t_j &\leq a_j+\delta+M(1-w_j).
  \label{eq:service}
\end{align}

Energy is propagated along each selected path. For direct rider and depot arcs:
\begin{align}
  b_j &\leq b_i-\epsilon(d_i,o_j)-\epsilon(o_j,d_j)+\bar B(1-z_{ij}), \\
  b_j &\leq b_i-\epsilon(d_i,\mathrm{dep})-\epsilon(\mathrm{dep},o_j)
        -\epsilon(o_j,d_j)+kg+2\bar B(1-x_{ijk}),
  \label{eq:energy}
\end{align}
where $g$ is the energy added per charging increment. Total charging is limited by the aggregate
capacity of the $n_C$ declared charge points:
\begin{equation}
  \textstyle\sum_{(i,j,k)}k x_{ijk}
  \leq n_C\left\lfloor\frac{T_1-T_0}{\kappa}\right\rfloor.
  \label{eq:charging}
\end{equation}

Additional valid inequalities limit total vehicle-minutes over the complete reporting window and over
selected subintervals. The objective is
\begin{equation}
  \max \textstyle\sum_{j\in\mathcal R} w_j.
  \label{eq:mip-objective}
\end{equation}


\section{Conditions and verification}
\label{sec:conditions}

\subsection{Attainable lower value}

The conservative solution records the assigned vehicle, pickup time, departure time, service mode and
request order for every selected rider. To establish $L$, the fleet follows that rider schedule inside
the reporting window and follows the originating policy outside it, under the same operating rules.
The latter condition reproduces the fleet state from which the benchmark instance was constructed.

The relation $L\leq N^*$ requires execution to reproduce the scheduled rider outcomes. Reconciliation
is performed rider by rider rather than only on the aggregate objective, because offsetting errors
could leave the total unchanged.

\subsection{Conditional upper value}

The relation $N^*\leq U$ requires every trajectory in $\Sigma$ to be representable in the optimistic
formulation. Three rules in the formulation are stricter than the declared operating rules and can
exclude feasible trajectories:
\begin{enumerate}
  \item a vehicle that waits in a request zone must arrive one complete invocation interval before the
        request;
  \item pickups after $T_1$ are excluded, although a request arriving before $T_1$ may be collected
        after the reporting window and still count as timely under the evaluation rules; and
  \item each rider receives at most one quote.
\end{enumerate}
Until their joint effect is measured, $U$ is a conditional upper comparison rather than a certified
upper bound on the unrestricted $N^*$.

Each solve also reports two diagnostic counts: selected riders whose vehicle departs after the
assignment deadline, and selected riders whose passenger journey duration does not match the travel
matrix at the planned pickup time. Both counts must be zero.

\end{document}
